\chapter{The 16-Week Implementation Roadmap}
\label{ch:roadmap}

The overall program from the remote-jobs playbook is framed as a 16-week process across seven phases. The source
material does not provide a fully specified week-by-week schedule within each phase, but it does establish the
phases themselves and what belongs in each. Where helpful, each phase below is cross-referenced to the chapter in
this guide that covers its detail.

\begin{longtable}{p{0.1\textwidth} p{0.24\textwidth} p{0.56\textwidth}}
\toprule
\textbf{Phase} & \textbf{Focus} & \textbf{What it involves} \\
\midrule
\endhead
1 & Build your foundation & Refine your technical positioning around judgment and business value
(Chapter~\ref{ch:foundations}). Identify the type of engineering work you want to be known for. Start a project
in which AI assists with implementation while you make and explain the important technical decisions
(Chapter~\ref{ch:projects}). \\
2 & Make your work production-oriented & Deploy the project, test it, monitor its behaviour, examine failure
cases, and document your decisions (Chapter~\ref{ch:projects-b}). Prepare demonstrations and technical
explanations that make your contribution clear (Chapter~\ref{ch:projects-c}). \\
3 & Establish your public presence & Update your professional profile and begin publishing meaningful technical
content (Chapter~\ref{ch:linkedin}). Share progress, decisions, problems, failures, and measured improvements.
Avoid generic or obviously automated posts. \\
4 & Build a target list & Identify funded startups and relevant founders. Research their products, business
models, recent launches, and likely technical needs. Select companies for which you have a credible value
proposition (Chapter~\ref{ch:direct-outreach}, Step 1). \\
5 & Start direct outreach & Begin contacting founders with tailored messages. Lead with a relevant observation or
proposed contribution rather than a generic request for employment. Track replies and learn from the responses
(Chapter~\ref{ch:direct-outreach}, Steps 2--4; the mechanics in Chapter~\ref{ch:outreach-engineering}). \\
6 & Rapid prototyping & For promising opportunities, develop small, relevant prototypes where appropriate.
Demonstrate what you built, how you evaluated it, and what limitations remain. Make clear that a prototype is a
demonstration, not an unsupported promise of production readiness (Chapter~\ref{ch:direct-outreach}, reciprocity
principle). \\
7 & Improve, follow up, and convert & Continue outreach, refine messaging, publish evidence of your work, and
build on promising conversations. If one campaign doesn't lead to an opportunity, learn from it and iterate. \\
\bottomrule
\end{longtable}

\begin{keytakeaway}
The hardest part is identified explicitly in the source material as \textbf{consistency}: maintaining this
process for 16 weeks while managing an existing full-time job or coursework. The difficulty is not purely
technical --- candidates must learn to write compelling messages, publish useful content, research companies,
select projects, and demonstrate professional judgment, all at once, over a sustained period.
\end{keytakeaway}

\section{A Suggested Weekly Rhythm}

None of the source documents specify exact weekly time blocks, but combining the 16-week phase structure with the
daily/weekly cadences given elsewhere in this guide (25--30 outreach messages/day once Phase 5 starts, 2+
LinkedIn posts/week from Phase 3 onward, 10+ substantive comments/day) suggests a realistic weekly rhythm once
you're past Phase 2:

\begin{itemize}
  \item \textbf{A fixed daily block} (even 30--45 minutes) for outreach and engagement --- this is what keeps $N$
  from collapsing to zero, per the Outreach Equation.
  \item \textbf{2--3 dedicated hours/week} for project work --- building, fixing, extending, or documenting the
  production-grade project from Part~III.
  \item \textbf{1--2 hours/week} for content --- drafting and posting from the content calendar in
  \S\ref{sec:linkedin-content}.
  \item \textbf{A weekly review} --- update your outreach CRM, check connection-acceptance and reply rates against
  the 40\%/20\% targets from Chapter~\ref{ch:linkedin}, and adjust messaging or targeting before the next week.
\end{itemize}

